\chapter{Background and Related Work}
\label{ch:background}

\section{Vision transformers}

\paragraph{ViT.} The vision transformer (ViT)~\cite{dosovitskiy2021vit} applies the encoder of the
transformer~\cite{vaswani2017attention} to images. An image of $224\times224$ pixels is cut into patches of $16\times16$ pixels. Each patch
is mapped to a vector (a \emph{token}), and a class token is added, which gives 197 tokens. The tokens pass through a
sequence of identical \emph{blocks}. Each block has two parts, each preceded by a layer normalisation and wrapped in
a residual connection:
\begin{itemize}
  \item \emph{self-attention}: three weight matrices (the \emph{query}, \emph{key} and \emph{value} projections) map
        every token to a query, a key and a value; every query is compared with every key, and the values are mixed
        accordingly; a fourth matrix (the output projection) follows;
  \item \emph{the MLP}: two linear layers with a GELU non-linearity between them. The first layer (fc1) widens the
        token to four times its width, the second (fc2) brings it back.
\end{itemize}
ViT-B/16, the model used here, has 12 blocks of width 768 and about 86 million weights. ViT-L has 24 blocks of
width 1024. DeiT-B~\cite{touvron2021deit} has exactly the architecture of ViT-B but was pretrained differently (on
ImageNet-1k only, with a stronger training recipe). The DeiT paper also introduces a variant with an additional
distillation token; the checkpoint we use is the one without it, so its inputs and blocks are those of ViT-B.

\paragraph{Swin.} The Swin transformer~\cite{liu2021swin} is hierarchical. Swin-B starts from $4\times4$ patches,
i.e.\ $56\times56 = 3136$ tokens of width 128, and works in four \emph{stages}. Between stages, neighbouring tokens
are merged: the number of tokens falls by a factor of four and the width doubles. Attention is computed only within
local $7\times7$ windows, which are shifted in every second block. Swin-B has 24 blocks in stages of 2, 2, 18 and 2
blocks (Table~\ref{tab:swin_stages}) and about 87 million weights. In this report its blocks are numbered 0 to 23
across the stages.

\begin{table}[H]
\centering\small
\begin{tabular}{lcccc}
\toprule
stage & blocks & tokens & width & block input per image (MiB) \\
\midrule
1 & 0--1 & $56\times56 = 3136$ & 128 & 1.53 \\
2 & 2--3 & $28\times28 = 784$ & 256 & 0.77 \\
3 & 4--21 & $14\times14 = 196$ & 512 & 0.38 \\
4 & 22--23 & $7\times7 = 49$ & 1024 & 0.19 \\
\midrule
ViT-B, all blocks & 0--11 & 197 & 768 & 0.58 \\
\bottomrule
\end{tabular}
\caption{The four stages of Swin-B, with ViT-B for comparison. Block input per image = tokens $\times$ width
$\times$ 4 bytes: the size of the tensor that enters a block, for one image in 32-bit floating point. The first
blocks of Swin-B work on the largest tensors.}
\label{tab:swin_stages}
\end{table}

\section{Low-rank adaptation}

LoRA~\cite{hu2022lora} adapts a linear layer with weight matrix $W \in \mathbb{R}^{d_\text{out}\times d_\text{in}}$
without changing $W$. It adds two small matrices, $A \in \mathbb{R}^{r\times d_\text{in}}$ and
$B \in \mathbb{R}^{d_\text{out}\times r}$, where the \emph{rank} $r$ is much smaller than the dimensions of $W$:
\begin{equation}
  y = Wx + \frac{\alpha}{r}\,B A x .
  \label{eq:lora}
\end{equation}
Only $A$ and $B$ (together the \emph{adapter}) are trained. $B$ starts at zero, so training starts from the
pretrained model exactly; $\alpha$ is a fixed scaling factor. With $r=8$ and width 768, an adapter has
$2 \times 8 \times 768 = 12{,}288$ values against 589,824 in $W$. Figure~\ref{fig:lora} shows the arrangement.

\begin{figure}[H]
\centering
\begin{tikzpicture}[font=\small, >={Stealth[length=2mm]},
  box/.style={draw, minimum height=0.9cm, minimum width=1.5cm, inner sep=3pt}]
  \node (x) at (0,0) {$x$};
  \node[box, fill=cGrey!25] (W) at (3.2,0.9) {$W$ (frozen)};
  \node[box, draw=cAttn, line width=1pt, fill=cAttn!12] (A) at (2.3,-0.9) {$A$};
  \node[box, draw=cAttn, line width=1pt, fill=cAttn!12] (B) at (4.3,-0.9) {$B$};
  \node[draw, circle, inner sep=1.5pt] (p) at (6.4,0) {$+$};
  \node (y) at (7.8,0) {$y$};
  \draw[->] (x) |- (W);
  \draw[->] (x) |- (A);
  \draw[->] (A) -- (B);
  \node[font=\scriptsize] at (3.3,-0.22) {rank $r$};
  \draw[->] (W) -| (p);
  \draw[->] (B) -| node[pos=0.25, below, font=\scriptsize] {$\times\,\alpha/r$} (p);
  \draw[->] (p) -- (y);
  \node[font=\scriptsize, text=cAttn, below=0.05cm of A] {trained};
  \node[font=\scriptsize, text=cAttn, below=0.05cm of B] {trained, starts at 0};
\end{tikzpicture}

\caption{A LoRA adapter on one linear layer. The pretrained weight $W$ is frozen; the low-rank path through $A$ and
$B$ is trained. After training, the product $\frac{\alpha}{r}BA$ can be added into $W$ (merging), which removes
the extra path.}
\label{fig:lora}
\end{figure}

\paragraph{Merging.} Because Equation~\ref{eq:lora} is linear in $x$, the adapter can be folded into the weight
after training: $W' = W + \frac{\alpha}{r}BA$. The merged layer is an ordinary linear layer, and the merged model
has the same structure and number of weights as the original one. We call a model without adapters, which is also
the form of a fully fine-tuned model, a \emph{dense} model. Subtracting the product again restores $W$
(\emph{unmerging}).

\paragraph{Where adapters can be placed.} Any linear layer can receive an adapter. In a transformer block the
natural choices are the attention projections and the two MLP layers; we call these the \emph{targets}. A second
choice is the set of blocks that receive adapters, which we call the \emph{placement}.

\section{What uses memory during training}
\label{sec:bg_memory}

The results of Chapters~\ref{ch:lora_ft} and~\ref{ch:placement} follow from one fact about backpropagation, so we
state it here. The memory needed for one training step has two parts.
\begin{itemize}
  \item \emph{Static memory}: the weights, one gradient per trainable weight, and the state of the optimiser (AdamW
        \cite{loshchilov2019adamw} keeps two more values per trainable weight). It does not depend on the batch.
        Full fine-tuning of an 86-million-weight model needs four values per weight (1.28~GiB); LoRA needs the
        weights once plus a few megabytes (0.33~GiB).
  \item \emph{Activation memory}: intermediate results of the forward pass that the backward pass will need. It
        grows in proportion to the batch size and, at the batch sizes used for training, it is by far the larger
        part: about 6.9 of the 8.2~GiB of full fine-tuning of ViT-B at batch size 64.
\end{itemize}
That stored activations, not weights, dominate the memory of training is the central observation of
TinyTL~\cite{cai2020tinytl}, which made it for convolutional networks.

What has to be stored depends on the operation. Let $x$ be the input of an operation, $y$ its output and
$\delta_y$ the gradient of the loss with respect to $y$, which arrives from the later layers.
\begin{itemize}
  \item A \emph{frozen linear layer} $y = Wx$ passes the gradient back as $\delta_x = W^\top \delta_y$. This needs only
        $W$. The input $x$ would be needed only for the gradient of $W$ itself, which is not computed. A frozen
        linear layer therefore stores nothing.
  \item A \emph{trainable layer}, such as a LoRA adapter, needs its input: the gradient of $A$ in
        Equation~\ref{eq:lora} is proportional to $(B^\top\delta_y)\,x^\top$. An adapter therefore forces the input
        of its layer to be stored.
  \item \emph{Layer normalisation, attention and GELU} are not multiplications by a fixed matrix. To pass a gradient
        back through them their inputs (and for attention also the output) must be stored, whether or not anything
        in them is trained.
  \item \emph{Blocks before the first trainable weight} are not part of the backward pass at all, and nothing in
        them is stored.
\end{itemize}
Two consequences follow, and both are measured in this report. First, the memory that LoRA saves depends on
\emph{which} layers have adapters: an adapter on fc2 forces the four-times-wider hidden layer of the MLP to be stored,
an adapter on an attention projection only the normal-width input of the block. Second, freezing a block does not
remove its cost as long as the backward pass still goes through it; the cost disappears only for blocks before the
first adapter (Figure~\ref{fig:reach}). The number of trainable weights says little about training memory.

\begin{figure}[H]
\centering
\begin{tikzpicture}[>={Stealth[length=2mm]},
  blk/.style={draw=cGrey, minimum width=0.82cm, minimum height=0.75cm, font=\scriptsize, inner sep=1pt},
  fwd/.style={blk, fill=cGrey!15, text=cGrey},
  ad/.style={blk, draw=cAttn, line width=1pt, fill=cAttn!20}]
  \node[fwd] (b0) at (0,0) {0};
  \foreach \i/\j in {1/0,2/1} {\node[fwd, right=0.1cm of b\j] (b\i) {\i};}
  \foreach \i/\j in {3/2,4/3,5/4,6/5,7/6,8/7,9/8,10/9,11/10} {\node[ad, right=0.1cm of b\j] (b\i) {\i};}
  \node[draw=cGrey, rounded corners=2pt, font=\scriptsize, right=0.25cm of b11, minimum height=0.75cm] (loss) {loss};
  \draw[->, cGrey, line width=0.8pt] ($(b0.north west)+(0,0.3)$) -- node[above, font=\scriptsize, text=black] {forward pass through all blocks} ($(loss.north east)+(0,0.3)$);
  \draw[->, cAccent, line width=1pt] ($(loss.south east)+(0,-0.3)$) -- node[below, font=\scriptsize, text=black] {backward pass: stops at the first adapted block} ($(b3.south west)+(0,-0.3)$);
  \node[font=\scriptsize, text=cGrey, align=center, below=0.35cm of b1] {nothing\\stored};
\end{tikzpicture}

\caption{Attention LoRA on blocks 3--11 of a 12-block model. The forward pass runs through all blocks. The backward
pass stops at the first adapted block, so only blocks 3--11 keep their intermediate results; blocks 0--2 store
nothing.}
\label{fig:reach}
\end{figure}

\section{Related work}
\label{sec:related}

This report is a measurement study, and we did not carry out a survey of the literature. Three works are directly
connected to it. LoRA~\cite{hu2022lora} is the method we study; it was proposed for language models, and we use its
implementation in the PEFT library~\cite{mangrulkar2022peft}. TinyTL~\cite{cai2020tinytl} observed, for
convolutional networks, that the memory of training is dominated by the stored activations and not by the number of
trained parameters; Chapters~\ref{ch:lora_ft} and~\ref{ch:placement} measure the same effect for LoRA in vision
transformers. Surgical fine-tuning~\cite{lee2023surgical} found that which layers are best to train depends on how
the new data differs from the pretraining data, which is related to our placement results
(Chapter~\ref{ch:placement}). We make no claim that the measurements reported here are new.
